% !TEX program = pdflatex
% Editable Beamer source. 16:9, no external figures, no .bib dependency.
% Replace \ShortSource footers or the final reference slides as you edit.
\documentclass[aspectratio=169,10pt]{beamer}

\usepackage[utf8]{inputenc}
\usepackage[T1]{fontenc}
\usepackage{lmodern}
\usepackage{booktabs}
\usepackage{tabularx}
\usepackage{array}
\usepackage{ragged2e}
\usepackage{xcolor}
\usepackage{hyperref}

\definecolor{Navy}{HTML}{17324D}
\definecolor{Blue}{HTML}{2F6690}
\definecolor{Teal}{HTML}{3A7D78}
\definecolor{Sand}{HTML}{F1E9DA}
\definecolor{Gold}{HTML}{D39B2A}
\definecolor{Brick}{HTML}{A84A3F}
\definecolor{Ink}{HTML}{20262E}
\definecolor{Muted}{HTML}{68727D}
\definecolor{Pale}{HTML}{F4F6F8}

\setbeamercolor{normal text}{fg=Ink,bg=white}
\setbeamercolor{structure}{fg=Navy}
\setbeamercolor{frametitle}{fg=white,bg=Navy}
\setbeamercolor{title}{fg=white}
\setbeamercolor{subtitle}{fg=Sand}
\setbeamercolor{block title}{fg=white,bg=Blue}
\setbeamercolor{block body}{fg=Ink,bg=Pale}
\setbeamercolor{alerted text}{fg=Brick}
\setbeamertemplate{navigation symbols}{}
\setbeamertemplate{itemize item}{\color{Teal}$\blacktriangleright$}
\setbeamertemplate{itemize subitem}{\color{Gold}$\bullet$}
\setbeamertemplate{footline}{%
  \hfill\usebeamerfont{page number in head/foot}\color{Muted}%
  \insertframenumber\hspace{0.7em}\vspace{0.45em}}
\setbeamersize{text margin left=0.62cm,text margin right=0.62cm}

\newcolumntype{Y}{>{\RaggedRight\arraybackslash}X}
\newcolumntype{L}[1]{>{\RaggedRight\arraybackslash}p{#1}}
\newcommand{\hi}{\textcolor{Teal}{\textbf{H}}}
\newcommand{\med}{\textcolor{Gold}{\textbf{M}}}
\newcommand{\low}{\textcolor{Brick}{\textbf{L}}}
\newcommand{\ShortSource}[1]{%
  \vfill\par\tiny\color{Muted}\textit{Sources:} #1}
\newcommand{\PolicyTag}[1]{%
  \colorbox{Sand}{\strut\textcolor{Navy}{\textbf{#1}}}}

\title{Four Brazilian policy experiments and internal migration}
\subtitle{Dams, FPM expansion, POLOCENTRO, and the BR-230/BR-163 frontier roads}
\author{Working deck for discussion and heavy revision}
\date{August 2026}

\begin{document}

{\setbeamercolor{background canvas}{bg=Navy}
\begin{frame}[plain]
  \vspace{1.15cm}
  {\usebeamerfont{title}\fontsize{25}{29}\selectfont\inserttitle\par}
  \vspace{0.35cm}
  {\usebeamerfont{subtitle}\Large\insertsubtitle\par}
  \vfill
  {\color{white}\small\insertauthor\par}
  {\color{Sand}\small\insertdate\par}
  \vspace{0.65cm}
\end{frame}}

% -----------------------------------------------------------------------------
% ONE-SLIDE CASE DESCRIPTIONS
% -----------------------------------------------------------------------------

\begin{frame}{1. Dam building: a staggered local boom plus forced displacement}
\small
\begin{columns}[T,onlytextwidth]
\begin{column}{0.49\textwidth}
  \PolicyTag{Policy shock}
  \begin{itemize}\setlength\itemsep{0.28em}
    \item Site-specific hydro projects bundle construction, reservoir flooding, resettlement, transmission, royalties, and compensatory spending.
    \item Treatment has several clocks: announcement/licensing, construction start, filling, operation, and royalty payments.
    \item Relevant geographies differ: host municipality, reservoir footprint, downstream communities, labor-sending places, and electricity-consuming regions.
  \end{itemize}
  \vspace{0.2em}
  \PolicyTag{Migration mechanisms}
  \begin{itemize}\setlength\itemsep{0.25em}
    \item Forced exit from inundated or livelihood-damaged areas.
    \item Temporary construction-worker inflow, followed by post-boom exit.
    \item Royalties/services can retain or attract residents; disrupted fishing/agriculture can expel them.
  \end{itemize}
\end{column}
\begin{column}{0.49\textwidth}
  \PolicyTag{What we already know}
  \begin{itemize}\setlength\itemsep{0.28em}
    \item \textbf{de Faria et al. (2017):} 56 plants, 1991--2010; GDP and tax booms were short-lived, with no detectable long-run social gains.
    \item \textbf{Costa, Szerman \& Assunção (2025):} 29 plants / 82 municipalities; average impacts trace an inverted U and show little sectoral transformation.
    \item \textbf{Randell (2016, 2017):} Belo Monte panel/interviews document forced migration, damaged livelihoods, and heterogeneous welfare effects.
  \end{itemize}
  \vspace{0.2em}
  \begin{block}{Research opening}
    Estimate \emph{who leaves, who arrives, and where they go}; then decompose place-based GDP/employment estimates into genuine productivity versus population composition and displacement.
  \end{block}
\end{column}
\end{columns}
\ShortSource{\href{https://doi.org/10.1016/j.eneco.2017.08.025}{de Faria et al. (2017)}; \href{https://doi.org/10.31389/eco.466}{Costa et al. (2025)}; \href{https://doi.org/10.1016/j.worlddev.2016.01.012}{Randell (2016)}.}
\end{frame}

\begin{frame}{2. FPM expansion: a fiscal amenity shock with endogenous population}
\small
\begin{columns}[T,onlytextwidth]
\begin{column}{0.49\textwidth}
  \PolicyTag{Policy shock}
  \begin{itemize}\setlength\itemsep{0.28em}
    \item The 1988 Constitution phased the municipal share of federal IR/IPI from \textbf{17\% to 22.5\%} by 1993, sharply expanding formula-based municipal revenue.
    \item For non-capital municipalities, allocations follow population brackets and coefficients codified in Decree-Law 1,881/1981.
    \item Two distinct experiments: \textbf{national expansion} (1988--93) versus \textbf{local bracket jumps} used in threshold RDD studies.
  \end{itemize}
  \vspace{0.2em}
  \PolicyTag{Migration mechanisms}
  \begin{itemize}\setlength\itemsep{0.25em}
    \item Public jobs, procurement, education/health, and local-demand multipliers may attract or retain residents.
    \item Capture, corruption, or low-salience spending can mute amenities; housing/land prices may capitalize gains.
  \end{itemize}
\end{column}
\begin{column}{0.49\textwidth}
  \PolicyTag{What we already know}
  \begin{itemize}\setlength\itemsep{0.28em}
    \item \textbf{Litschig \& Morrison (2013):} extra transfers raised spending, schooling and literacy, and reduced poverty.
    \item \textbf{Brollo et al. (2013):} more revenue increased observed corruption and reduced candidate education.
    \item \textbf{Corbi et al. (2019):} sizable transfer-driven local employment/output multiplier.
    \item Direct migration evidence is mainly conference/working-paper work and is mixed; one RDD study reports lower inflow and higher outflow.
  \end{itemize}
  \begin{block}{Research opening}
    Treat migration as both an outcome and a threat: the running variable is population itself. Separate the 1988 reform from the modern threshold design.
  \end{block}
\end{column}
\end{columns}
\ShortSource{\href{https://doi.org/10.1257/app.5.4.206}{Litschig \& Morrison (2013)}; \href{https://doi.org/10.1257/aer.103.5.1759}{Brollo et al. (2013)}; \href{https://doi.org/10.1093/restud/rdy064}{Corbi et al. (2019)}; \href{https://www.anpec.org.br/encontro/2019/submissao/files_I/i12-c352774646be9eef49d2a1d4a639819a.pdf}{ANPEC migration paper}.}
\end{frame}

\begin{frame}{3. POLOCENTRO: subsidized frontier modernization, 1975--early 1980s}
\small
\begin{columns}[T,onlytextwidth]
\begin{column}{0.49\textwidth}
  \PolicyTag{Policy shock}
  \begin{itemize}\setlength\itemsep{0.28em}
    \item Created by Decree 75,320/1975 to modernize agriculture in selected Cerrado poles in the Center-West and western Minas Gerais.
    \item Intended to incorporate roughly \textbf{3.7 million hectares} during 1975--79; retrospective sources identify about \textbf{202 municipalities}.
    \item A bundled treatment: subsidized rural credit, roads, electrification, storage, research, extension, and marketing support.
  \end{itemize}
  \vspace{0.2em}
  \PolicyTag{Migration mechanisms}
  \begin{itemize}\setlength\itemsep{0.25em}
    \item Attracted capitalized farmers and technicians, especially from the South/Southeast.
    \item Mechanization, land appreciation, and consolidation reduced labor demand and displaced smallholders.
    \item New farm/service towns absorbed some rural exit while other workers moved beyond the treated poles.
  \end{itemize}
\end{column}
\begin{column}{0.49\textwidth}
  \PolicyTag{What we already know}
  \begin{itemize}\setlength\itemsep{0.28em}
    \item \textbf{FJP (1984):} a nine-volume official evaluation documents implementation, production change, and strong concentration of credit/benefits.
    \item \textbf{Mueller (1990), Salim (1986), later geography/political economy:} faster commercial farming and soy expansion, rising land concentration, weak employment intensity, and small-producer marginalization.
    \item The literature is rich descriptively but contains little modern causal work and almost no direct migration evaluation.
  \end{itemize}
  \begin{block}{Research opening}
    Reconstruct pole boundaries and annual program intensity; distinguish selected producer in-migration from incumbent displacement and rural-to-urban exit.
  \end{block}
\end{column}
\end{columns}
\ShortSource{\href{https://www2.camara.leg.br/legin/fed/decret/1970-1979/decreto-75320-29-janeiro-1975-423871-publicacaooriginal-1-pe.html}{Decree 75,320/1975}; \href{https://www.bibliotecadigital.mg.gov.br/consulta/consultaDetalheDocumento.php?iCodDocumento=55206}{FJP (1984)}; \href{https://repositorio.ipea.gov.br/bitstream/11058/14764/1/ppp_n3_completo.pdf}{Mueller (1990)}.}
\end{frame}

\begin{frame}{4. BR-230 and BR-163: roads plus state-directed frontier settlement}
\small
\begin{columns}[T,onlytextwidth]
\begin{column}{0.49\textwidth}
  \PolicyTag{Policy shock}
  \begin{itemize}\setlength\itemsep{0.28em}
    \item \textbf{Transamazônica (BR-230):} construction began in 1970 under the PIN; a major segment opened in 1972 alongside INCRA colonization and planned agrovilas.
    \item \textbf{Cuiabá--Santarém (BR-163):} constructed from 1971 and inaugurated in 1976; later evolved from colonization axis to logging/cattle/soy corridor.
    \item Treatment is actual accessibility, not a line on a map: segment opening, paving, bridges, maintenance, feeder roads, and seasonal passability matter.
  \end{itemize}
  \vspace{0.2em}
  \PolicyTag{Migration mechanisms}
  \begin{itemize}\setlength\itemsep{0.25em}
    \item Lower moving/trade costs plus land grants directly recruited settlers.
    \item Induced logging, cattle, farming, land speculation, and road towns attracted secondary flows.
    \item Project failure and weak services generated turnover and onward/out-migration.
  \end{itemize}
\end{column}
\begin{column}{0.49\textwidth}
  \PolicyTag{What we already know}
  \begin{itemize}\setlength\itemsep{0.28em}
    \item Transamazon anthropology/geography documents low fulfillment of settlement targets, difficult soils/services, high turnover, and major social/environmental disruption.
    \item BR-163 studies document evolving mobility, land conflict, deforestation, and a shift toward export-oriented agribusiness.
    \item Causal economics for these two corridors is thin. Related Brazil evidence shows roads reduce both trade and migration costs and reshape population.
  \end{itemize}
  \begin{block}{Research opening}
    Estimate segment-level event studies with actual passability and separate road access from INCRA settlement, credit, and subsequent feeder-road expansion.
  \end{block}
\end{column}
\end{columns}
\ShortSource{\href{https://doi.org/10.1590/1809-4422ASOC256V1922016}{Bratman (2016)}; \href{https://doi.org/10.1525/9780520314320}{Smith (1982)}; \href{https://docs.rwu.edu/fcas_fp/151/}{Campbell (2012/14)}; \href{https://doi.org/10.1257/app.20180487}{Morten \& Oliveira (2024)}.}
\end{frame}

% -----------------------------------------------------------------------------
% COMPARISON TABLES
% -----------------------------------------------------------------------------

\begin{frame}{IPEA evaluability dimensions I: policy definition and exposure}
\scriptsize
\textit{Provisional mapping: the exact public ``Índice de Potencial Avaliativo'' scorecard was not located. Rows follow IPEA's public catalog/evaluability fields. H/M/L are this deck's qualitative ratings, not official IPEA scores.}
\vspace{0.35em}
\renewcommand{\arraystretch}{1.25}
\begin{tabularx}{\textwidth}{L{2.35cm}YYYY}
\toprule
\textbf{Dimension} & \textbf{Dams} & \textbf{FPM expansion} & \textbf{POLOCENTRO} & \textbf{BR-230 / BR-163} \\
\midrule
\textbf{Objective / theory of change} & \hi\ Clear output objective; local-development claims secondary and heterogeneous. & \hi\ Fiscal decentralization/equalization; spending discretion makes downstream theory broad. & \med\ Modernize/occupy selected Cerrado poles through credit and infrastructure. & \med\ Integration, settlement, security, and frontier development were bundled goals. \\
\textbf{Target population / eligibility} & \med\ Reservoir-affected people can be mapped, but ``affected'' extends downstream and through labor markets. & \hi\ Formula and bracket rules are explicit; capitals/reserve municipalities differ. & \med\ Pole/municipality lists exist; beneficiary rules and realized access were selective. & \med\ Corridor buffers and INCRA projects are recoverable; actual access varied by segment/season. \\
\textbf{Implementation / rollout} & \hi\ Project milestones are dated, although announcement, works, filling, and operation differ. & \hi\ Constitutional phase-in and annual coefficients are observable. & \med\ 1975 start is clear; annual rollout and closure differ by component/place. & \low--\med\ Opening dates are broad; completion, paving, maintenance, and passability require reconstruction. \\
\textbf{Exposure measurement} & \hi\ Coordinates, reservoir polygons, capacity, stage, royalties. & \hi\ Theoretical and actual transfers; pre-reform dependence. & \low--\med\ Treatment often binary in secondary sources; credit/infrastructure intensity fragmented. & \med\ Distance/network access feasible, but nominal road presence overstates usable access. \\
\bottomrule
\end{tabularx}
\ShortSource{IPEA's catalog lists objectives, description, target public/criteria, budget, results, monitoring reports, and evaluations as core attributes: \href{https://catalogo.ipea.gov.br/arquivos/posts/4565-td2799web.pdf}{Brito, Mello \& Alencar (2022), pp. 7--8}.}
\end{frame}

\begin{frame}{IPEA evaluability dimensions II: resources, outcomes, and evidence}
\scriptsize
\renewcommand{\arraystretch}{1.25}
\begin{tabularx}{\textwidth}{L{2.35cm}YYYY}
\toprule
\textbf{Dimension} & \textbf{Dams} & \textbf{FPM expansion} & \textbf{POLOCENTRO} & \textbf{BR-230 / BR-163} \\
\midrule
\textbf{Budget / treatment intensity} & \med--\hi\ Capex, capacity, inundation, and royalties exist but across agencies/projects. & \hi\ Treasury/TCU transfers and statutory formula are annual and municipal. & \low--\med\ Aggregate budgets known; municipal credit and component spending need archives. & \low--\med\ Construction/paving contracts are fragmented; physical accessibility may be the better dose. \\
\textbf{Monitoring / output records} & \hi\ ANEEL/EPE, licensing/EIA, resettlement records, satellite reservoirs. & \hi\ Treasury, TCU coefficients, municipal accounts, public-sector employment. & \med\ FJP evaluation and agency archives; limited standardized annual monitoring. & \med\ DNER/DNIT/INCRA archives plus Landsat; road condition is poorly standardized historically. \\
\textbf{Outcome data} & \hi\ Census migration, RAIS, GDP/taxes, royalties, health/violence, land use. & \hi\ Census migration, FINBRA/SICONFI, education/health, RAIS; population assignment needs care. & \med\ 1970/1980/1991 censuses and agricultural censuses; sparse annual pre-period. & \med--\hi\ Census/locality data, agricultural census, imagery, deforestation; exact historical access is the bottleneck. \\
\textbf{Previous evaluation} & \hi\ Published causal econ plus displacement case studies. & \hi\ Extensive published RDD literature; little published migration work. & \med\ Large official/descriptive literature; almost no modern causal identification. & \med\ Rich qualitative/environmental evidence; thin causal work on these exact corridors. \\
\textbf{Overall potential} & \hi\ Strong, especially with displacement microdata. & \hi\ Best clean first stage, but migration endogeneity is fundamental. & \med\ High novelty; data reconstruction is the project. & \med\ High substantive value; treatment reconstruction and bundling are difficult. \\
\bottomrule
\end{tabularx}
\end{frame}

\begin{frame}{What does the existing literature actually establish?}
\scriptsize
\renewcommand{\arraystretch}{1.22}
\begin{tabularx}{\textwidth}{L{2.05cm}L{3.35cm}YL{3.25cm}}
\toprule
\textbf{Experiment} & \textbf{Dominant literature type} & \textbf{Main findings} & \textbf{Migration evidence} \\
\midrule
\textbf{Dams} & Published applied/environmental econ; demography; mixed-method displacement; EIA/NGO studies. & Short construction boom in GDP, employment, lights, and taxes; little average long-run transformation; large heterogeneity and livelihood losses among displaced groups. & Direct, credible forced-displacement studies for Belo Monte; multi-dam causal work rarely follows origin--destination flows or decomposes population composition. \\
\textbf{FPM} & Published top-field/general-interest econ using RDD; public finance reports; recent working papers. & More spending and local activity; improved schooling/literacy and lower poverty in one historical design; more corruption/lower politician quality in another. & Thin and not settled: conference papers use census migration and thresholds, with some counterintuitive outflow results. \\
\textbf{POLOCENTRO} & Official FJP/IPEA evaluation; economic history; rural sociology; geography/political ecology. & Commercial agriculture/soy expansion, infrastructure and land appreciation; credit and land concentration; mechanization and weak employment absorption; smallholder marginalization. & Mostly inferred from demographic change, rural exodus, and case narratives; little person-level or causal origin--destination evidence. \\
\textbf{BR-230 / BR-163} & Anthropology/geography/history; World Bank/government reports; environmental science/remote sensing. Related causal transport econ. & Directed colonization underperformed official targets; high turnover and service failures; roads enabled settlement, deforestation, land conflict, cattle/soy/logging, and new towns. & Mobility central in qualitative work; little corridor-specific causal estimation. Brasília-road evidence validates trade- and migration-cost channels elsewhere in Brazil. \\
\bottomrule
\end{tabularx}
\end{frame}

\begin{frame}{Proposed mechanisms: why should each policy change migration?}
\scriptsize
\renewcommand{\arraystretch}{1.22}
\begin{tabularx}{\textwidth}{L{2.15cm}YYY}
\toprule
\textbf{Experiment} & \textbf{Pull / retention} & \textbf{Push / displacement} & \textbf{Empirical signature} \\
\midrule
\textbf{Dams} & Construction jobs; procurement; royalties; compensatory services; improved local access. & Inundation/resettlement; lost land/fisheries; congestion, prices, violence; post-construction job collapse. & Inflow rises near works, then reverses; forced movers cluster by resettlement policy; composition changes even if total population does not. \\
\textbf{FPM} & Public employment, services, schooling/health, private multiplier, better amenities. & Capture/corruption; tax or effort substitution; housing-price capitalization; skilled exit if services disappoint. & Higher inflow/lower outflow if amenities dominate; sector/skill heterogeneity; potential feedback across population brackets. \\
\textbf{POLOCENTRO} & Subsidized credit/infrastructure raises farm returns; attracts capitalized producers, technicians, traders, and urban services. & Mechanization lowers labor demand; land consolidation/appreciation displaces smallholders and tenants. & Simultaneous skilled/landed in-migration and incumbent rural out-migration; farm-size and occupation composition shift sharply. \\
\textbf{BR-230 / BR-163} & Lower moving/trade costs; land grants; resource jobs; market access; agglomeration and road towns. & Indigenous/peasant displacement; conflict; failed crops/services; seasonal isolation; frontier boom-bust. & Distance-decay around newly usable segments; strong origin networks; turnover and onward migration; feeder-road spillovers. \\
\bottomrule
\end{tabularx}
\end{frame}

\begin{frame}{Threats to causal identification}
\scriptsize
\renewcommand{\arraystretch}{1.22}
\begin{tabularx}{\textwidth}{L{2.05cm}YY}
\toprule
\textbf{Experiment} & \textbf{Selection / confounding} & \textbf{Measurement / estimand problems} \\
\midrule
\textbf{Dams} & Endogenous siting in high-potential basins; concurrent roads/transmission; anticipation and litigation; project cancellation is non-random. & Multiple treatment dates; reservoir/downstream/grid spillovers; administrative boundaries; municipality averages omit displaced people and destination effects. \\
\textbf{FPM} & 1988 reform coincides with democratization/decentralization; threshold-adjacent policies; manipulation of official population; endogenous municipal creation. & Population is both running variable and migration outcome; imperfect formula compliance; local RDD estimand; cross-border fiscal and labor spillovers. \\
\textbf{POLOCENTRO} & Poles selected for infrastructure/agronomic potential; concurrent EMBRAPA technology, credit, prices, roads, and state programs; beneficiary selection. & Fuzzy/changing boundaries; little annual intensity data; few pre-periods; municipality splits; program spillovers and producer sorting. \\
\textbf{BR-230 / BR-163} & Routes, sequencing, paving, and maintenance target strategic/economic places; INCRA land, credit, and services co-treated; endogenous feeder roads. & Nominal opening $\neq$ usable access; segment dates uncertain; wide spatial spillovers; anticipation/speculation; corridor outcomes mix road and colonization effects. \\
\bottomrule
\end{tabularx}
\end{frame}

\begin{frame}{Design responses: feasible ways to reduce those threats}
\scriptsize
\renewcommand{\arraystretch}{1.19}
\begin{tabularx}{\textwidth}{L{2.05cm}YL{4.0cm}}
\toprule
\textbf{Experiment} & \textbf{Core design} & \textbf{Diagnostics and extensions} \\
\midrule
\textbf{Dams} & Event study / synthetic control using planned-but-unbuilt or not-yet-built projects; date announcement, works, filling, operation separately; map reservoir/downstream rings. & Pre-trends and placebo dates; heterogeneity by inundation/royalties/capacity; track census microdata origins and destinations; aggregate origin + destination welfare; control co-built roads. \\
\textbf{FPM} & Fuzzy RDD: theoretical coefficient instruments actual transfers. For 1988 expansion, interact predetermined municipal FPM dependence with the national phase-in schedule. & Density/donut tests; predetermined census population; omit overlapping thresholds and new municipalities; estimate inflow and outflow separately; assess neighbor spillovers and local nature of RDD. \\
\textbf{POLOCENTRO} & Digitize official pole boundaries and annual credit/infrastructure; boundary DiD or matched/synthetic controls using 1970 baseline; dose-response by realized program intensity. & Build stable 1970 AMCs; document co-programs; test pre-treatment farm structure; heterogeneity by land size/soil suitability; use agricultural and population censuses jointly. \\
\textbf{BR-230 / BR-163} & Segment-level staggered event study using actual opening/passability; planned-but-unbuilt segments or historical engineering cost as counterfactual; network market-access treatment. & Separate road buffers from INCRA projects; remote sensing/archival validation; ring and spillover estimates; stable AMCs/localities; placebo future segments; model origin networks and feeder-road growth. \\
\bottomrule
\end{tabularx}
\end{frame}

\begin{frame}{Research portfolio: identification, novelty, and migration payoff}
\small
\renewcommand{\arraystretch}{1.28}
\begin{tabularx}{\textwidth}{L{2.5cm}L{2.2cm}L{2.2cm}L{2.2cm}Y}
\toprule
 & \textbf{First stage} & \textbf{Data burden} & \textbf{Novelty} & \textbf{Best contribution} \\
\midrule
\textbf{Dams} & Strong / staggered & Medium & High & Put out-migration and displaced people back into local-development estimates. \\
\textbf{FPM} & Very strong at cutoffs; weaker for 1988 reform & Low--medium & Medium & Clean fiscal-amenity test, but migration feedback to population is the central design issue. \\
\textbf{POLOCENTRO} & Medium / bundled & High & Very high & Causal account of agricultural modernization as simultaneous migrant selection and displacement. \\
\textbf{BR-230 / BR-163} & Medium if segment timing is recovered & Very high & Very high & Quantify state-built migration corridors and distinguish access from colonization policy. \\
\bottomrule
\end{tabularx}
\vspace{0.55em}
\begin{block}{Suggested sequencing}
\textbf{Fastest credible paper:} FPM threshold design. \quad
\textbf{Best fit to the substantive question:} dams. \quad
\textbf{Highest archival upside:} POLOCENTRO or the two Amazon corridors.
\end{block}
\begin{alertblock}{Common outcome architecture}
Estimate in-migration, out-migration, net migration, migrant skill/occupation, and origin--destination flows separately. Net population change alone can hide large offsetting flows and selection.
\end{alertblock}
\end{frame}

% -----------------------------------------------------------------------------
% REFERENCES (kept inside the .tex for easy editing)
% -----------------------------------------------------------------------------

\begin{frame}{Selected references I: evaluability, dams, and FPM}
\tiny
\begin{itemize}\setlength\itemsep{0.35em}
  \item Brito, A.; Mello, J.; Alencar, J. (2022). \href{https://catalogo.ipea.gov.br/arquivos/posts/4565-td2799web.pdf}{\textit{Catálogo de Políticas Públicas: primeiros resultados e hipóteses de pesquisa}. IPEA TD 2799.}
  \item IPEA / Casa Civil (2018). \href{https://ipea.gov.br/portal/images/stories/PDFs/livros/livros/181218_avaliacao_de_politicas_publicas_vol2_guia_expost.pdf}{\textit{Avaliação de Políticas Públicas: Guia Prático de Análise Ex Post}, vol. 2.}
  \item de Faria, F. A. M.; Davis, A.; Severnini, E.; Jaramillo, P. (2017). \href{https://doi.org/10.1016/j.eneco.2017.08.025}{The local socio-economic impacts of large hydropower plant development in a developing country. \textit{Energy Economics} 67: 533--544.}
  \item Costa, F.; Szerman, D.; Assunção, J. (2025). \href{https://doi.org/10.31389/eco.466}{Local Economic Impacts of Hydroelectric Power Plants: Evidence from Brazil. \textit{Economía}.}
  \item Randell, H. (2016). \href{https://doi.org/10.1016/j.worlddev.2016.01.012}{The short-term impacts of development-induced displacement on wealth and subjective well-being in the Brazilian Amazon. \textit{World Development} 87.}
  \item Randell, H. (2017). \href{https://pmc.ncbi.nlm.nih.gov/articles/PMC5347396/}{Structure and agency in development-induced forced migration: the case of Brazil's Belo Monte Dam. \textit{Population and Environment}.}
  \item Litschig, S.; Morrison, K. (2013). \href{https://doi.org/10.1257/app.5.4.206}{The Impact of Intergovernmental Transfers on Education Outcomes and Poverty Reduction. \textit{AEJ: Applied Economics} 5(4): 206--240.}
  \item Brollo, F.; Nannicini, T.; Perotti, R.; Tabellini, G. (2013). \href{https://doi.org/10.1257/aer.103.5.1759}{The Political Resource Curse. \textit{American Economic Review} 103(5): 1759--1796.}
  \item Corbi, R.; Papaioannou, E.; Surico, P. (2019). \href{https://doi.org/10.1093/restud/rdy064}{Regional Transfer Multipliers. \textit{Review of Economic Studies} 86(5): 1901--1934.}
  \item Monasterio, L. (2013). \href{https://doi.org/10.1016/j.jpubeco.2012.11.001}{Are rules-based government programs shielded from special-interest politics? \textit{Journal of Public Economics} 106.}
  \item Mendes, M.; Miranda, R. B.; Cossio, F. B. (2008). \href{https://www12.senado.leg.br/publicacoes/estudos-legislativos/tipos-de-estudos/outras-publicacoes/volume-iv-constituicao-de-1988-o-brasil-20-anos-depois.-estado-e-economia-em-vinte-anos-de-mudancas/do-sistema-tributario-nacional-o-fundo-de-participacao-dos-municipios-precisa-mudar}{O Fundo de Participação dos Municípios precisa mudar. Senado Federal.}
\end{itemize}
\end{frame}

\begin{frame}{Selected references II: POLOCENTRO and frontier roads}
\tiny
\begin{itemize}\setlength\itemsep{0.35em}
  \item Brasil (1975). \href{https://www2.camara.leg.br/legin/fed/decret/1970-1979/decreto-75320-29-janeiro-1975-423871-publicacaooriginal-1-pe.html}{Decreto 75.320, de 29 de janeiro de 1975 (POLOCENTRO).}
  \item Fundação João Pinheiro (1984). \href{https://www.bibliotecadigital.mg.gov.br/consulta/consultaDetalheDocumento.php?iCodDocumento=55206}{\textit{Estudos para a redefinição do Programa de Desenvolvimento dos Cerrados -- POLOCENTRO}, 9 vols.}
  \item Mueller, C. C. (1990). \href{https://repositorio.ipea.gov.br/bitstream/11058/14764/1/ppp_n3_completo.pdf}{Políticas governamentais e a expansão recente da agropecuária no Centro-Oeste. \textit{Planejamento e Políticas Públicas} 3: 45--74.}
  \item Salim, C. A. (1986). \href{https://doi.org/10.35977/0104-1096.cct1986.v3.9213}{As políticas econômica e tecnológica para o desenvolvimento agrário das áreas de cerrados no Brasil. \textit{Cadernos de Ciência \& Tecnologia} 3.}
  \item Smith, N. J. H. (1982). \href{https://doi.org/10.1525/9780520314320}{\textit{Rainforest Corridors: The Transamazon Colonization Scheme}. University of California Press.}
  \item Bratman, E. (2016). \href{https://doi.org/10.1590/1809-4422ASOC256V1922016}{Roads and dams: infrastructure-driven transformations in the Brazilian Amazon. \textit{Ambiente \& Sociedade} 19(2).}
  \item Campbell, J. M. (2012/2014). \href{https://docs.rwu.edu/fcas_fp/151/}{Brazil's Deferred Highway: Mobility, Development, and Anticipating the State in Amazonia. \textit{Boletín de Antropología} 27(44).}
  \item Fearnside, P. M. (2007). \href{https://repositorio.inpa.gov.br/items/72737248-5a0c-42dd-adf3-82b4a9d622d6}{Brazil's Cuiabá--Santarém (BR-163) Highway: the environmental cost of paving a soybean corridor through the Amazon. \textit{Environmental Management}.}
  \item Soares-Filho, B. et al. (2004). \href{https://ipam.org.br/bibliotecas/simulating-the-response-of-land-cover-changes-to-road-paving-and-governance-along-a-major-amazon-highway-the-santarem-cuiaba-corridor/}{Simulating land-cover changes to road paving and governance along the Santarém--Cuiabá corridor. \textit{Global Change Biology}.}
  \item Morten, M.; Oliveira, J. (2024). \href{https://doi.org/10.1257/app.20180487}{The Effects of Roads on Trade and Migration: Evidence from a Planned Capital City. \textit{AEJ: Applied Economics} 16(2): 389--421.}
  \item World Bank (various years). Project and environmental reviews of Amazon settlement and road programs; useful for implementation facts, not treated here as causal estimates.
\end{itemize}
\vfill
\tiny\color{Muted}\textit{Editing note:} The deck intentionally keeps citations inline and uses no BibTeX file. Replace the provisional IPEA checklist with the exact IPA rubric if you have the internal scoring sheet.
\end{frame}

\end{document}
